\documentclass[11pt,a4paper]{article}
\newcommand{\fichetype}{Exercices}
% Thème A4 commun, librement modifiable.
\usepackage[T1]{fontenc}
\usepackage[utf8]{inputenc}
\usepackage[french]{babel}
\usepackage{lmodern,amsmath,amssymb,mathtools,tikz,pgfplots}
\usepackage[margin=2.2cm,headheight=15pt]{geometry}
\usepackage{xcolor,fancyhdr,enumitem,hyperref,needspace}
\usetikzlibrary{arrows.meta,positioning,calc}
\pgfplotsset{compat=1.18,width=11cm,height=6cm}
\definecolor{Ink}{HTML}{202124}
\definecolor{Blue}{HTML}{356B9E}
\definecolor{Muted}{HTML}{737579}
\color{Ink}
\renewcommand{\familydefault}{\sfdefault}
\setlength{\parindent}{0pt}
\setlength{\parskip}{6pt}
\setlist{itemsep=3pt,topsep=4pt}
\pagestyle{fancy}\fancyhf{}
\fancyhead[L]{\small\color{Muted}MATHS LYCÉE}
\fancyhead[R]{\small\color{Muted}2026--2027}
\fancyfoot[L]{\small\color{Muted}\fichetype}
\fancyfoot[R]{\small\thepage}
\renewcommand{\headrulewidth}{0.3pt}
\hypersetup{colorlinks=true,urlcolor=Blue,linkcolor=Blue}
\newcommand{\R}{\mathbb{R}}
\newcommand{\N}{\mathbb{N}}
\newcommand{\E}{\mathbb{E}}
\newcommand{\V}{\mathrm{V}}
\newcommand{\exercice}[3]{\Needspace{6\baselineskip}\vspace{0.5cm}{\color{Blue}\small\textbf{EXERCICE #1}\hfill #2}\par{\large\bfseries #3}\par}

\begin{document}
{\small\color{Blue}PREMIÈRE SPÉCIALITÉ}\par{\LARGE\bfseries Nombre dérivé, tangente et approximation}\par{\large Exercices}\par\vspace{0.2cm}\hrule\vspace{0.4cm}
Les exercices progressent de l’application directe vers la recherche. Détailler les calculs et justifier les réponses.
\par\noindent\begin{minipage}{\linewidth}\exercice{1}{Application}{Taux moyens}
Pour $f(x)=2x^2+1$, calculer les taux entre $1$ et $3$, puis entre $1$ et $1,1$. Interpréter géométriquement.
\end{minipage}\par
\par\noindent\begin{minipage}{\linewidth}\exercice{2}{Entraînement}{Une dérivée par définition}
Pour $f(x)=3x^2-2x$, calculer le taux en $a$ avec l’accroissement $h\ne0$, puis $f\prime(a)$.
\end{minipage}\par
\par\noindent\begin{minipage}{\linewidth}\exercice{3}{Application}{Équation de tangente}
$f(x)=x^2-4x+1$. On admet $f\prime(x)=2x-4$. Déterminer la tangente en $a=3$ et son intersection avec l’axe des ordonnées.
\end{minipage}\par
\par\noindent\begin{minipage}{\linewidth}\exercice{4}{Application}{Lecture d’une tangente}
La courbe passe par $A(-1,4)$. Sa tangente en $A$ passe par $B(2,-2)$. Donner $f(-1)$, $f\prime(-1)$ et l’équation de cette tangente.
\end{minipage}\par
\par\noindent\begin{minipage}{\linewidth}\exercice{5}{Entraînement}{Approximer un carré}
Utiliser la tangente à la fonction carré en $5$ pour approcher $5,03^2$. Calculer ensuite l’erreur exacte et son signe.
\end{minipage}\par
\par\noindent\begin{minipage}{\linewidth}\exercice{6}{Entraînement}{Une vitesse instantanée}
Une distance fictive vaut $d(t)=2t^2$ mètres. Calculer la vitesse moyenne entre $3$ et $3+h$, puis la vitesse instantanée en $3$.
\end{minipage}\par
\par\noindent\begin{minipage}{\linewidth}\exercice{7}{Entraînement}{Une liste de sécantes}
Pour $f(x)=1/x$, calculer les taux en $1$ pour $h=0,1$ et $h=-0,1$. Décrire un algorithme produisant les taux pour $h=10^{-k}$, $k=1,\ldots,5$.
\end{minipage}\par
\par\noindent\begin{minipage}{\linewidth}\exercice{8}{Approfondissement}{Deux fonctions en zéro}
Comparer la dérivabilité en $0$ de $f(x)=|x|$ et $g(x)=x|x|$. Utiliser les quotients de variation.
\end{minipage}\par
\end{document}
